\documentclass{article}
\usepackage{amsmath,amssymb,hyperref}
\begin{document}
\section*{Jacobson's conjecture for two-sided Noetherian rings}
\subsection{Definitions and mathematical statement}
Let $R$ be an associative ring with identity. Assume that ascending chains of
left ideals stabilize and that ascending chains of right ideals stabilize.
Write $J(R)$ for the intersection of its maximal left ideals, the Jacobson
radical. For $n\geq1$, the ideal $J(R)^n$ consists of finite sums of products
of $n$ radical elements.

Must every such ring satisfy
\[
 \bigcap_{n\geq1}J(R)^n=\{0\}?
\]
Thus an element belonging to every positive power would have to vanish;
the question does not require a single power of the radical to be zero.
Both Noetherian assumptions are part of the question.

\subsection{Short English statement}
In a ring that is Noetherian on both sides, must the powers of the Jacobson
radical have zero intersection?

\subsection{Sources}
N. Jacobson, \emph{Structure of Rings} (1956), p.\ 200, for the historical
one-sided question. The two-sided formulation and the distinction from known
one-sided counterexamples are reviewed in A. Reyes,
\emph{Jacobson's conjecture and skew PBW extensions}, Revista Integraci\'on
32 (2014), 139--152, especially p.\ 140,
\url{https://dialnet.unirioja.es/descarga/articulo/5752599.pdf}.

A. V. Jategaonkar, \emph{Jacobson's conjecture and modules over fully bounded
Noetherian rings}, Journal of Algebra 30 (1974), 103--121,
\url{https://doi.org/10.1016/0021-8693(74)90195-1}, proves an important
special case. The statement above is an independently written formulation
of the cited mathematical question; source exposition has not been reproduced.
\end{document}
